\documentclass[10pt,a4paper]{article}

% --- UNIVERSAL PREAMBLE BLOCK ---
\usepackage[a4paper, top=3cm, bottom=2.5cm, left=2.2cm, right=2cm, headheight=35pt, headsep=12pt]{geometry}
\usepackage[T1]{fontenc}
\usepackage[utf8]{inputenc}
\usepackage{lmodern}
\usepackage[english]{babel}

% PDFLaTeX-compatible font setup

\usepackage{enumitem}
\setlist[itemize]{label=-}
\setlist[enumerate]{itemsep=0.15em,topsep=0.2em}

% --- REQUIRED PACKAGES ---
\usepackage{tabularx}
\usepackage{longtable}
\usepackage{tcolorbox}
\usepackage{fancyhdr}
\usepackage{lastpage}
\usepackage{ragged2e}
\usepackage{array}
\usepackage{url}

% --- SETTINGS & COMMANDS ---
\setlength{\parindent}{0pt}
\setlength{\parskip}{0.5em}
\renewcommand{\arraystretch}{1.3}

% Tcolorbox setup for mimicking form boxes
\tcbset{
    colback=white,
    colframe=black,
    sharp corners,
    boxrule=0.5pt,
    boxsep=2pt,
    left=6pt,
    right=6pt,
    top=6pt,
    bottom=6pt
}

% Custom commands
\newcommand{\sectionlabel}[1]{\vspace{1.2em}\noindent\textbf{\fontsize{10.5}{12}\selectfont #1}\par\vspace{0.4em}}
\newcommand{\subhead}[1]{\vspace{0.8em}\noindent\textbf{\fontsize{10}{11}\selectfont #1}\par\vspace{0.3em}}
\newcommand{\formline}{\vspace{1em}\noindent\rule{\textwidth}{0.5pt}\vspace{1em}}

% Column types for tables
\newcolumntype{Y}{>{\RaggedRight\arraybackslash}X}
\newcolumntype{C}{>{\centering\arraybackslash}X}
\newcolumntype{P}[1]{>{\RaggedRight\arraybackslash}p{#1}}

% --- HEADER & FOOTER SETTINGS ---
\pagestyle{fancy}
\fancyhf{}
\fancyhead[C]{\fontsize{9}{10}\selectfont\bfseries 2209/A UNIVERSITY STUDENTS RESEARCH PROJECTS SUPPORT PROGRAM\\RESEARCH PROPOSAL FORM}
\fancyfoot[C]{\fontsize{9}{10}\selectfont \thepage~/~\pageref{LastPage}}
\renewcommand{\headrulewidth}{0pt}

\fancypagestyle{firstpage}{
    \fancyhf{}
    \fancyfoot[C]{\fontsize{9}{10}\selectfont \thepage~/~\pageref{LastPage}}
    \renewcommand{\headrulewidth}{0pt}
}

\begin{document}
\thispagestyle{firstpage}

% --- COVER PAGE ---
\vspace*{3cm}
\begin{center}
    \textbf{\Huge TÜBİTAK}\\[2.5cm]
    {\Large\bfseries TÜBİTAK-2209-A UNIVERSITY STUDENTS RESEARCH PROJECTS SUPPORT PROGRAM}\\[1.5cm]
    {\large\bfseries RESEARCH PROPOSAL FORM}\\[2.5cm]
\end{center}
\vfill
\begin{flushright}
    \textbf{[Yıl] Year -- [Dönem] Period Application}
\end{flushright}

\clearpage

% --- PAGE 2 ONWARDS ---

\sectionlabel{A. GENERAL INFORMATION}
\begin{tcolorbox}
\textbf{Applicant's Full Name:} [Öğrenci Adı Soyadı] \\[0.6em]
\textbf{Title of the Research Proposal:} [Proje Başlığı] \\[0.6em]
\textbf{Advisor's Full Name:} [Danışman Adı Soyadı] \\[0.6em]
\textbf{Institution/Organization Where the Research Will Be Conducted:} [Üniversite/Kurum Adı]
\end{tcolorbox}

\sectionlabel{ABSTRACT}
\begin{tcolorbox}
{[Projenin Türkçe veya İngilizce özeti buraya yazılacaktır. Konunun genel çerçevesi, projenin amacı, yöntemi ve beklenen sonuçları kısaca özetlenmelidir...]}

\vspace{2em}
\textbf{Keywords:} [Anahtar Kelime 1], [Anahtar Kelime 2], [Anahtar Kelime 3]
\end{tcolorbox}

\sectionlabel{ORIGINAL VALUE}
\subhead{Importance of the Subject, Original Value of the Research Proposal and Research Question/Hypothesis}
\textbf{1.1. Importance of the Subject and Scientific Excellence of the Project}

AYBUStore addresses a structural access problem arising from the university's multi-campus organization. The service is currently physically available at only one location, while students are distributed across four campuses. This configuration creates unequal access to basic daily supplies, increases time and effort costs, and limits service continuity for students who are geographically distant from the store.

The project's scope is the digital transformation of this single-point physical service into a university-centered, hyper-local e-commerce system. Its boundary is intra-university use; it is designed specifically for the internal dynamics of the AYBU student ecosystem rather than for national-scale commerce. In this framework, the project is expected to increase service accessibility by at least 30\% and to improve platform responsiveness through measurable Time to Interactive (TTI) optimization.

From a scientific perspective, the project targets a gap not sufficiently addressed by conventional e-commerce studies and platforms: the design of an equitable, closed-loop digital supply model for a specific multi-campus university context. Unlike general-market systems, AYBUStore proposes a context-specific architecture that integrates local demand patterns, campus distribution constraints, and continuous digital access requirements into a single service model.

The original contribution is fourfold. Conceptually, the project introduces a ``Hyper-Local Closed-Loop Digital Ecosystem'' for intra-university commerce. Methodologically, it replaces the traditional physical supply process with an inclusive and digitized service delivery model suitable for distributed campuses. Technically, it combines a modern front-end, scalable back-end, and embedded automated support modules in an end-to-end architecture. Practically, it enables equal, rapid, and location-independent access for students across all campuses on a 24/7 basis, thereby strengthening service equity and institutional integration.

\textbf{1.2. Research Question and/or Hypothesis}

Within a multi-campus university structure, a single-point physical market model produces unequal access, time loss, and operational inefficiency in meeting students' routine needs. This project investigates whether a hyper-local, closed-loop digital commerce model can eliminate these limitations.

The study addresses the following questions:
\begin{enumerate}
    \item How can a university-specific hyper-local digital supply model be designed to provide continuous and equitable cross-campus access?
    \item To what extent can a component-based, scalable architecture improve system responsiveness (TTI) relative to current static portal structures?
    \item How effectively can integrated support modules improve support speed and operational resolution performance?
\end{enumerate}

The hypotheses are:
\begin{itemize}
    \item \textbf{H1 (Inclusivity and Access):} AYBUStore will increase cross-campus service access by at least 30\% and reduce spatial inequality compared with the current physical model.
    \item \textbf{H2 (Technical Performance):} The proposed modern architecture will improve Time to Interactive (TTI) by at least 40\% versus the baseline static web portal and reduce average TTI below 1.5 seconds.
    \item \textbf{H3 (Operational Efficiency):} Integrated rule-based support modules (asynchronous chatbots, automated ticket routing, dynamic notifications) will reduce initial support response time by 60\% and average issue resolution time by at least 45\% relative to manual processes.
\end{itemize}

Hypotheses will be tested through a 3-month pilot after system completion. Log-based indicators (TTI, order process time, ticket response and resolution times) from at least 300 active students ($n \geq 300$), selected via stratified sampling, will be compared with historical baseline data from the existing system. Statistical significance will be evaluated using Independent Samples t-test (or Mann-Whitney U when distributional assumptions are not met), with significance threshold set at $p < 0.05$.

\subhead{Aims and Objectives}
\textbf{1.3. Aim and Objectives of the Project}

The main aim of this project is to develop AYBUStore as a hyper-local, closed-loop digital commerce ecosystem that removes physical access constraints in a multi-campus university setting and provides equitable, location-independent, and continuous service access.

To achieve this aim, the project defines the following measurable objectives for the 3-month pilot period:
\begin{itemize}
    \item \textbf{Objective 1 (Access and Inclusivity):} Increase peripheral-campus utilization by at least 30\% relative to the baseline (current estimated active access: 25\% of total students) and reach at least 300 active verified student users across campuses.
    \item \textbf{Objective 2 (Technical Performance):} Reduce average Time to Interactive (TTI) from the baseline of approximately 3.5 seconds to below 1.5 seconds (minimum 40\% improvement).
    \item \textbf{Objective 3 (Operational Efficiency):} Reduce first support response time from 12 hours by 60\% and average issue-resolution time from 48 hours by at least 45\% through integrated rule-based support modules.
    \item \textbf{Objective 4 (Reliability and Scalability):} Maintain 99.9\% uptime and 0\% transaction-level data loss during the pilot while sustaining peak load conditions of 500 requests/second and 1,000 concurrent active sessions.
\end{itemize}

Measurement is defined as follows: an active user is a unique SIS-verified student performing at least one in-system action within a 30-day window; utilization rate is calculated as
\[
\textit{Utilization Rate (\%)} = \left(\frac{U_{\mathrm{active,\,peripheral}}}{U_{\mathrm{total,\,peripheral}}}\right) \times 100
\]
where $U_{\mathrm{active,\,peripheral}}$ denotes active peripheral-campus users and $U_{\mathrm{total,\,peripheral}}$ denotes total peripheral-campus students.

Reliability and performance will be monitored via continuous APM-based server logging (e.g., Prometheus/Grafana), and data integrity will be verified through PostgreSQL transaction/WAL logs and automated ACID compliance tests.

\subhead{METHOD}
[Araştırmada kullanılacak yöntem, veri toplama teknikleri, kullanılacak araçlar ve analiz yöntemleri bu bölümde detaylandırılacaktır...]

\clearpage

\sectionlabel{PROJECT MANAGEMENT}
\subhead{Work-Time Schedule}
\begin{center}\textbf{WORK-TIME SCHEDULE}\end{center}
\begin{tabularx}{\textwidth}{|c|P{0.25\textwidth}|Y|P{0.12\textwidth}|Y|}
\hline
\textbf{IP No} & \textbf{Name and Objectives of Work Packages} & \textbf{By whom(s) it will be carried out} & \textbf{Time Interval (Month)} & \textbf{Success Criteria and Contribution to the Success of the Project} \\ \hline
1 & [İş Paketi 1 Adı ve Amacı] & [Sorumlu Kişi(ler)] & Month [X-Y] & [Başarı Kriteri 1] \\ \hline
2 & [İş Paketi 2 Adı ve Amacı] & [Sorumlu Kişi(ler)] & Month [X-Y] & [Başarı Kriteri 2] \\ \hline
3 & [İş Paketi 3 Adı ve Amacı] & [Sorumlu Kişi(ler)] & Month [X-Y] & [Başarı Kriteri 3] \\ \hline
\end{tabularx}

\vspace{1.5em}
\subhead{Risk Management}
\begin{center}\textbf{RISK MANAGEMENT TABLE*}\end{center}
\begin{tabularx}{\textwidth}{|c|Y|Y|}
\hline
\textbf{IP No} & \textbf{Most Significant Risks} & \textbf{Risk Management (Plan B)} \\ \hline
1 & [Olası Risk 1] & [B Planı / Çözüm Önerisi 1] \\ \hline
2 & [Olası Risk 2] & [B Planı / Çözüm Önerisi 2] \\ \hline
3 & [Olası Risk 3] & [B Planı / Çözüm Önerisi 3] \\ \hline
\end{tabularx}\\[0.5em]
\small{(*) The rows in the table can be expanded and reproduced as necessary.}

\vspace{1.5em}
\subhead{3.3. Research Opportunities}
\begin{center}\textbf{TABLE OF RESEARCH OPPORTUNITIES (*)}\end{center}
\begin{tabularx}{\textwidth}{|Y|Y|}
\hline
\textbf{Type and Model of Infrastructure/Equipment in the Organization (Laboratory, Vehicle, Machinery-Equipment, etc.)} & \textbf{Intended Use in the Project} \\ \hline
[Altyapı/Ekipman 1] & [Projedeki Kullanım Amacı 1] \\ \hline
[Altyapı/Ekipman 2] & [Projedeki Kullanım Amacı 2] \\ \hline
[Altyapı/Ekipman 3] & [Projedeki Kullanım Amacı 3] \\ \hline
\end{tabularx}\\[0.5em]
\small{(*) The rows in the table can be expanded and reproduced as necessary.}


\clearpage
\sectionlabel{4. IMPACT}

If the project is carried out successfully, AYBUStore is expected to produce scientific, technical, educational and social outputs by transforming a single-location university store model into a measurable, campus-wide digital commerce ecosystem. The project will contribute to equal access to university-related products and services across distributed campuses, improve operational visibility for administrative units, and provide a reusable software engineering case for smart campus digital transformation.

\subhead{4.1. Expected Outputs}

\small
\setlength\LTleft{0pt}
\setlength\LTright{0pt}
\begin{longtable}{|P{0.27\textwidth}|P{0.46\textwidth}|P{0.20\textwidth}|}
\hline
\textbf{Output type} & \textbf{Expected Output(s)} & \textbf{Expected Time Interval} \\ \hline
\textbf{Scientific/Academic Outputs}\newline {\small (National/International Article, Book, Book chapter, Proceeding etc.)}
& The project will produce a structured academic case study on the design of a \textbf{hyper-local closed-loop digital commerce ecosystem} for a multi-campus university. Expected outputs include a final technical research report, system architecture documentation, pilot evaluation results and a student conference paper or article draft. The pilot data will provide measurable evidence on cross-campus access improvement, Time to Interactive performance and support-process efficiency.
& \textbf{0--12 months:} final technical report, system design documentation and pilot evaluation outputs.\par \textbf{12--18 months:} student conference paper or article draft. \\ \hline
\textbf{Economical/}\newline\textbf{Commercial/}\newline\textbf{Social Outputs}\newline {\small (Prototype, Product, Patent, Utility Model, Production License, Inventory/Database/Documentation Production, Spin-off/Start-up Company etc.)}
& The main output will be the \textbf{AYBUStore prototype/reference implementation}, including OTP-based authentication, product catalogue, cart/order flow, simulated checkout, stock management, admin/CMS panel, campus-aware tracking and rule-based support modules. The project will also produce database/inventory structures for users, products, categories, orders, support tickets and tracking data, together with SRS, UML diagrams, test cases and user/admin documentation. The expected social output is a digital service model that reduces access inequality caused by the current single-location store structure.
& \textbf{0--12 months:} prototype/reference implementation, database model, SRS, UML diagrams, test cases and user/admin documentation.\par \textbf{Post-project:} possible institutional adaptation by university units or extension into a broader smart-campus service model. \\ \hline
\textbf{Outputs Regarding to}\newline\textbf{Training Researcher and}\newline\textbf{Creating New Project(s)}\newline {\small (Graduate thesis, National/International new project)}
& The project will train undergraduate researchers in requirements engineering, Agile/Scrum planning, software architecture, cloud-based backend design, secure OTP authentication, database modelling, testing and performance evaluation. It will also provide experience in hypothesis-based research through pilot data, log-based measurement and statistical comparison. The outputs can support future extensions or new projects on smart campus services, AI-assisted support systems, campus-aware order tracking, demand analysis and digital transformation in higher education.
& \textbf{0--12 months:} undergraduate researcher training and reusable project artefacts.\par \textbf{12--18 months / post-project:} possible new student projects, graduation projects or institutional digital transformation initiatives. \\ \hline
\end{longtable}
\normalsize

{\small (*) Time intervals are specified as 0--12 months, 12--18 months and post-project where applicable.}

\subhead{4.2. Expected Impacts}

\textbf{Expected Application Areas:}

The primary application area of AYBUStore is \textbf{university-centered digital commerce and smart campus service delivery}. The project is designed for Ankara Yıldırım Beyazıt University's distributed campus structure, especially Esenboğa, Etlik, Bilkent and Çubuk campuses. The system can be used as a digital platform for campus-focused product discovery, order management, inventory visibility, support handling and campus-aware order tracking.

The potential end users of the project outputs are students, especially those located on peripheral campuses; university administrative units and campus store operators; university IT units; and software engineering students or researchers. Students will benefit from easier access to university-branded products, academic supplies and campus-related merchandise. Administrative and store units will benefit from better product, stock, order and support-process visibility. IT units may use the architecture and documentation as a reference for future smart-campus services. Researchers and students may use the project as a case study for requirements engineering, cloud-based architecture, secure authentication, usability evaluation and performance measurement.

In the short term, the project outputs are expected to be applied within the AYBU campus ecosystem as a prototype/reference implementation. In the medium and long term, the same model may be adapted by other universities that operate across multiple campuses and face similar problems related to service accessibility, inventory management and student-centered digital transformation.

\textbf{Socio-economic/Cultural Contribution:}

AYBUStore is expected to contribute to \textbf{social equity and service accessibility} within the university by reducing the access gap between students located on different campuses. In the current model, students who are physically distant from the main store location may face additional time, transportation and opportunity costs. By digitizing product discovery, order placement, tracking and support processes, the project supports more equal participation in campus life.

From a socio-economic perspective, the project may reduce unnecessary physical travel for routine product access and improve operational efficiency through digital stock tracking, order monitoring and administrative control. This can help university units make more informed decisions about product availability, campus-based demand and service continuity.

From an educational perspective, the project contributes to the development of qualified undergraduate researchers in software engineering. The team will gain practical experience in project management, requirements analysis, system design, secure authentication, cloud services, testing, risk management and technical documentation.

The project is also aligned with broader digital transformation and smart campus goals. It supports student-centered service design, improvement of education-related service quality, inclusive access to institutional services and the use of scalable digital infrastructure in higher education.

\subhead{4.3. Dissemination of Project Results and Activity Plan to be Implemented within the Scope of Science Communication}

\small
\begin{longtable}{|P{0.24\textwidth}|P{0.69\textwidth}|}
\hline
\textbf{Field} & \textbf{Content} \\ \hline
\textbf{Target Group} & The target groups are AYBU students, academic and administrative staff, university IT units, campus store administrators, software engineering students, researchers working on digital transformation in higher education, and other universities interested in smart-campus service models. Students will benefit from improved access to campus-related products and services. Administrative and IT units will benefit from the system architecture, database model, documentation and operational workflow. Software engineering students and researchers will benefit from the project as a real-world case study. \\ \hline
\textbf{Goals and Expected Gains} & The main goal is to increase awareness of how a campus-specific digital commerce system can improve cross-campus accessibility, operational efficiency and student experience. Expected gains include better understanding of service inequality in multi-campus universities, increased institutional awareness of smart campus solutions, and creation of reusable technical knowledge for future student projects or institutional digital transformation initiatives. Dissemination activities will also support feedback collection from potential users and stakeholders. \\ \hline
\textbf{Tools to be Used} & The project results will be shared through a final project presentation, technical report, prototype demonstration, UML diagrams, system architecture documentation, screenshots, short demo video and project poster/infographic. If institutional permission is obtained, selected outputs may also be shared through department seminars, student project exhibitions, university digital platforms or a controlled GitHub/code archive. These tools are selected because they are suitable for both technical and non-technical target groups. \\ \hline
\textbf{Timing} & Dissemination will begin during the final project presentation period within \textbf{0--12 months}. The technical report, prototype demonstration, documentation and evaluation results will be shared at the end of the development cycle. Feedback-oriented meetings or demonstrations with relevant university stakeholders may be conducted after the prototype/reference implementation is completed. Academic dissemination such as a student conference paper or article draft may be prepared within \textbf{12--18 months}. \\ \hline
\end{longtable}
\normalsize

\sectionlabel{ANY OTHER ISSUES YOU WOULD LIKE TO MENTION}
[Buraya projeyle ilgili eklemek istediğiniz diğer detaylar, etik kurallar, izin belgeleri veya destek alınan farklı hususlar eklenecektir...]

\sectionlabel{APPENDICES ANNEX-1: RESOURCES}
\formline
\begin{enumerate}[leftmargin=1.4em]
\item Alenezi, M., Akour, M. 2023. ``Digital transformation blueprint in higher education: A case study of PSU'', \textit{Sustainability}, 15 (10), 8204. DOI: 10.3390/su15108204.
\item Alsheyadi, A. K., Albalushi, J. 2020. ``Service quality of student services and student satisfaction: The mediating effect of cross-functional collaboration'', \textit{The TQM Journal}. DOI: 10.1108/TQM-10-2019-0234.
\item Brdesee, H., Alsaggaf, W. 2022. ``Decision-making strategy for digital transformation: A two-year analytical study and follow-up concerning innovative improvements in university e-services'', \textit{Journal of Theoretical and Applied Electronic Commerce Research}, 17 (1), 138--164. DOI: 10.3390/jtaer17010008.
\item Dinh, H., Tran, T. K. 2023. ``EduChat: An AI-based chatbot for university-related information using a hybrid approach'', \textit{Applied Sciences}, 13 (22), 12446. DOI: 10.3390/app132212446.
\item Gong, T., Yi, Y. 2018. ``The effect of service quality on customer satisfaction, loyalty, and happiness in five Asian countries'', \textit{Psychology \& Marketing}, 35 (6), 427--442. DOI: 10.1002/mar.21096.
\item Güldal, H., Dinçer, E. O. 2025. ``Can rule-based educational chatbots be an acceptable alternative for students in higher education?'', \textit{Education and Information Technologies}, 30, 3979--4012. DOI: 10.1007/s10639-024-12977-5.
\item Jahromi, H. Z., Delaney, D. T., Hines, A. 2020. ``Beyond first impressions: Estimating quality of experience for interactive web applications'', \textit{IEEE Access}, 8, 47741--47755. DOI: 10.1109/ACCESS.2020.2979385.
\item Liu, D., Deng, Z., Zhang, W., Wang, Y., Kaisar, E. I. 2021. ``Design of sustainable urban electronic grocery distribution network'', \textit{Alexandria Engineering Journal}, 60 (1), 145--157. DOI: 10.1016/j.aej.2020.06.051.
\item Melkonyan, A., Gruchmann, T., Lohmar, F., Bleischwitz, R., Spinler, S. 2020. ``Sustainability assessment of last-mile logistics and distribution strategies: The case of local food networks'', \textit{International Journal of Production Economics}, 228, 107746. DOI: 10.1016/j.ijpe.2020.107746.
\item Netravali, R., Nathan, V., Mickens, J., Balakrishnan, H. 2018. ``Vesper: Measuring time-to-interactivity for web pages'', \textit{15th USENIX Symposium on Networked Systems Design and Implementation (NSDI 18)}, 217--231. İnternet adresi: \url{https://www.usenix.org/conference/nsdi18/presentation/netravali-vesper}, Son erişim tarihi: 12 Mayıs 2026.
\item Suryawanshi, P., Dutta, P., L, V., G, D. 2021. ``Sustainable and resilience planning for the supply chain of online hyperlocal grocery services'', \textit{Sustainable Production and Consumption}, 28, 496--518. DOI: 10.1016/j.spc.2021.05.001.
\item Urquhart, R., Newing, A., Hood, N., Heppenstall, A. 2022. ``Last-mile capacity constraints in online grocery fulfilment in Great Britain'', \textit{Journal of Theoretical and Applied Electronic Commerce Research}, 17 (2), 636--651. DOI: 10.3390/jtaer17020033.
\end{enumerate}
\end{document}
